\section{验证：Shell 命令}
\label{sec:vma-verify}

\subsection{\texttt{vma list}——列出所有 VMA}

\begin{tcblisting}{consolebox}
«\tpk{>}» «\tin{vma list}»
=== VMA list (sorted-array, 2 entries) ===
  [0xC0000000, 0xD0000000) rwx  direct-map      (256 MiB)
  [0xE0000000, 0xE0010000) rw-  kheap           (64 KiB)
\end{tcblisting}

初始状态有两个 VMA：直接映射区和内核堆。地址按升序排列（红黑树中序遍历），
权限用 \texttt{rwx} 紧凑表示。

\subsection{\texttt{vma find <addr>}——查找指定地址}

\begin{tcblisting}{consolebox}
«\tpk{>}» «\tin{vma find 0xC1000000}»
0xC1000000 -> VMA [0xC0000000, 0xD0000000) 'direct-map'
  size: 262144 KiB, offset: 0x1000000, flags: rwx

«\tpk{>}» «\tin{vma find 0xA0000000}»
0xA0000000 -> no VMA (unmapped / untracked)
\end{tcblisting}

\subsection{\texttt{vma test}——自测试}

\begin{tcblisting}{consolebox}
«\tpk{>}» «\tin{vma test}»
[vma-test] === VMA selftest ===
[vma-test] test 1: add + find
[vma-test] PASS
[vma-test] test 2: find outside
[vma-test] PASS
[vma-test] test 3: overlap detection
[vma-test] PASS
[vma-test] test 4: remove + find
[vma-test] PASS
[vma-test] === ALL PASS ===
\end{tcblisting}

测试使用 \hex{F1000000}--\hex{F1010000}（64\,KiB）作为临时区域，
验证 add/find/overlap/remove 四个基本操作的正确性。

% ============================================================================

